\subsubsection{General Inflow Boundary Condition}
\label{SEC:BC:Inflow:General}

This boundary condition combines subsonic and supersonic inflow boundary
conditions into a single option.
%
This boundary condition takes the velocity from the left\slash interior state,
$\vec{u}_\subL$, and evaluates the magnitude of that velocity normal to the
boundary face, $v_{n,\subL}$, as
%
\[
  v_{n,\subL} = \vec{u}_\subL\cdot\normal
\]
%
where \normal{} is the unit normal for the boundary face that points outside the
computational domain.
%
The interior static temperature $T_\subL$ and $v_{n,\subL}$ can then be used to
compute an interior Mach number normal to the inflow boundary, $M_{n,\subL}$, as
%
\[
  M_{n,\subL} = \frac{v_{n,\subL}}{\sqrt{\gamma \rgas T_\subL}}
\]
%
If $M_{n,\subL} \leq -1$ (negative because the normal points outside the
domain), the inflow is determined to be supersonic, and the right state is
evaluated the same way as the supersonic inflow of
\sref{SEC:BC:Inflow:Supersonic} by just fixing it to
the boundary condition state, \thatis \[ \cvVecR = \cvVecBC \]
%
Otherwise, the inflow is subsonic, and the right state is evaluated based on the
subsonic inflow boundary condition with $\inptext{Sub\_Inflow\_Method}=4$
(\ref{SEC:BC:Inflow:Subsonic:4}).

%%%%%%%%%%%%%%%%%%%%%%%%%%%%%%%%%%%%%%%%%%%%%%%%%%%%%%%%%%%%%%%%%%%%%%%%%%%%%%%%

It should be noted that this boundary condition will always operate as if
$\inptext{Sub\_Inflow\_Method}=4$ when the interior Mach number is found to be
subsonic.  The value of the input \inptext{Sub\_Inflow\_Method} has no control
over the subsonic inflow method used for this boundary condition and is
therefore ignored.

%%%%%%%%%%%%%%%%%%%%%%%%%%%%%%%%%%%%%%%%%%%%%%%%%%%%%%%%%%%%%%%%%%%%%%%%%%%%%%%%

\devNotes{Need to fix subsonic inflow angle-of-attack bug.}

%###############################################################################

